\chapter*{Vorwort}
\addcontentsline{toc}{chapter}{Vorwort}

Im Rahmen meines Verbundstudiums bei der infoteam Software AG konnte ich bereits vor Beginn meiner Bachelorarbeit in der Abteilung „Life Science“ wertvolle Einblicke in die Entwicklung von Komponenten für zenLAB gewinnen.
Dabei mein besonderes Interesse an der Weiterentwicklung dieser Softwarelösung entstanden. 
Da die Anbindung von Geräten in der Medizintechnik eine zentrale und zugleich anspruchsvolle Aufgabe darstellt, beschäftigt sich diese Bachelorarbeit mit der Erweiterung der bestehenden Software 
um ein Device-Integration-Plugin. 
Ziel der Arbeit ist es, eine technische Grundlage zu schaffen, durch die neue Geräte künftig einfacher, flexibler und strukturierter in zenLAB integriert werden können.

An dieser Stelle möchte ich mich herzlich bei Jonas Krammer und Dr. Melanie Kahl für die Betreuung im Unternehmen, 
die fachliche Unterstützung sowie das hilfreiche Feedback während der Erstellung meiner Bachelorarbeit bedanken. 
Mein Dank gilt außerdem dem gesamten zenLAB Team für die interessanten Gespräche und den konstruktiven Austausch zur Entwicklung und Integration des Device-Integration-Plugins. 
Insbesondere durch gemeinsame Diskussionen konnten unterschiedliche Perspektiven eingebracht, Herausforderungen frühzeitig erkannt und mögliche Lösungsansätze weiterentwickelt werden.

Darüber hinaus danke ich Fabio Starke für seine Überarbeitungsvorschläge zur schriftlichen Ausarbeitung. 
Ebenso möchte ich mich bei Prof. Dr. Matthias Teßmann von der Technischen Hochschule Nürnberg für die Betreuung meiner Bachelorarbeit, 
die Beantwortung meiner Fragen sowie das wertvolle Feedback bedanken.

In der vorliegenden Arbeit werden Eigenbegriffe, 
Markennamen von Unternehmen oder staatlichen Einrichtungen sowie Begriffe aus dem Glossar bei ihrer ersten Nennung kursiv hervorgehoben. 
Codeausschnitte, beispielsweise Klassen- oder Methodennamen, werden im Fließtext farblich gekennzeichnet. 
Aus Gründen der besseren Lesbarkeit wird in dieser Arbeit das generische Maskulinum verwendet. 
Dieses ist geschlechtsneutral zu verstehen und bezieht sich gleichermaßen auf Personen jeglichen Geschlechts.
